% TOP_PROBLEM: {"id": 30006801, "title": "Berry random-wave conjecture: real negative-curvature BS formulation", "category_id": 16, "category": {"id": 16, "name": "physics", "display_name": "Mathematical Physics", "description": "Problems at the intersection of mathematics and physics.", "slug": "physics", "order_index": 16, "created_at": "2026-07-31T15:26:25.671Z"}, "status": "open"}
% FIRST_POSTED: 2026-09-25
% ENTRY_KIND: statement_only
% !TeX program = lualatex
% Definition and sources only; this file is not an LLM solution attempt.
\documentclass[11pt]{article}
\usepackage[margin=25mm]{geometry}
\usepackage{fontspec}
\setmainfont{Latin Modern Roman}
\usepackage{amsmath,amssymb,mathtools}
\usepackage{unicode-math}
\setmathfont{Latin Modern Math}
\usepackage{xurl,hyperref}
\hypersetup{colorlinks=true,urlcolor=blue}
\setcounter{secnumdepth}{0}
\setlength{\parindent}{0pt}
\setlength{\parskip}{5pt}
\setlength{\emergencystretch}{3em}
\begin{document}
\section*{\texorpdfstring{Berry random-wave conjecture: real negative-curvature BS formulation}{Berry random-wave conjecture: real negative-curvature BS formulation}}\label{30006801}
\subsection{Definitions and mathematical statement}
Let $\mu=\operatorname{vol}_g/\operatorname{vol}_g(M)$ on a closed connected negatively curved $d$-manifold, and choose any real orthonormal eigenbasis with
$\Delta_g\phi_j=\alpha_j^2\phi_j$, $\alpha_j>0$, ordered by eigenvalue. Sample $p$ from $\mu$. The target is weak convergence of the laws of the pointed decorated manifolds
\[
 [M,\alpha_j^2g,p,\phi_j]\ \Longrightarrow\ [{\mathbb{R}}^d,g_{\mathrm E},0,F]
\]
in the source's local smooth topology on pointed isometry classes. The limiting real Gaussian field has mean zero and covariance
${\mathbb{E}}[F(x)F(y)]=\int_{S^{d-1}}\cos\langle x-y,\theta\rangle{\,\mathrm{d}}\sigma(\theta)$,
where $\sigma$ is probability surface measure. Thus it has unit variance and unit spectral radius. Convergence is along the full positive-energy sequence for every real eigenbasis, not just a density-one subsequence or a single moment test.
\par\smallskip\textit{Formulation and concept references:} \href{https://arxiv.org/pdf/1810.05601v2}{[S1]}.

\subsection{Short English statement}
When viewed at wavelength scale around a random point, must every high-energy eigenfunction on a negatively curved manifold statistically resemble the universal real Gaussian random wave?
\subsection{Sources}
\begin{itemize}
\item[S1]Eigenfunctions and random waves in the Benjamini-Schramm limit; arXiv1810.05601v2.
\url{https://arxiv.org/pdf/1810.05601v2}.
\textit{Formulation source.}
\end{itemize}
\end{document}
